% Vietnamese translation for OI Public Library
% Source-File: IOI中国国家候选队论文/2005/杨思雨/试题/登山问题.doc
\documentclass[11pt]{article}
\usepackage{oipl-vn}

\title{Bài toán leo núi}
\author{SHTSC 2002}
\date{}

\begin{document}
\maketitle

\origtitle{EXPLORER}
\sourcepath{IOI Candidate Team Papers/2005/Yang Siyu/problems/EXPLORER.doc}

\section{Phát biểu}

Một nhóm OIer chuẩn bị vượt qua OI Mountain, một dãy núi hùng vĩ. Trước khi
bắt đầu, họ muốn biết vị trí đỉnh cao nhất của dãy núi.

Mọi mặt cắt ngang của dãy núi là như nhau. Ở hai phía của đỉnh cao nhất, độ
cao giảm dần theo thứ tự. Các OIer có thể dùng hệ thống vệ tinh để đo độ cao
của núi tại một vị trí cụ thể. Mỗi lần phát lệnh \texttt{Explorer(x)}, vệ
tinh trả về độ cao của điểm có tọa độ \(x\) trên mặt cắt ngang.

Do việc liên lạc với vệ tinh luôn chậm, họ muốn tìm được vị trí đỉnh cao nhất
trong số lượt hỏi rất ít. Bạn có làm được không?

\section{Thư viện tương tác}

Thư viện kiểm tra cung cấp ba hàm: \texttt{Start}, \texttt{Ask} và
\texttt{Answer}.
\begin{itemize}
  \item \texttt{Start()} phải được gọi đầu tiên; hàm này trả về giá trị thực
  \(L\).
  \item \texttt{Ask(X)} dùng để hỏi giá trị \(f(X)\), tức độ cao tại vị trí
  \(X\).
  \item \texttt{Answer(Ans)} dùng để báo cho thư viện kiểm tra rằng chương
  trình cho rằng \(f(x)\) đạt cực đại tại \(x=\texttt{Ans}\).
\end{itemize}
Chương trình không được tự kết thúc và không được truy cập bất kỳ tệp nào.

\section{Giao diện cho Pascal}

Chương trình Pascal cần dùng câu lệnh sau để nạp thư viện:
\begin{verbatim}
uses tools_p;
\end{verbatim}

Các khai báo do thư viện cung cấp là:
\begin{verbatim}
function Start: double;
function Ask(x: double): double;
procedure Answer(ans: double);
\end{verbatim}

\section{Giao diện cho C++}

Với C++, thêm dòng sau vào đầu chương trình:
\begin{verbatim}
#include "tools_c.h"
\end{verbatim}
Trong RHIDE, đặt tùy chọn \texttt{Option->Linker} là \texttt{tools\_c.o}.

Các nguyên mẫu hàm là:
\begin{verbatim}
double Start();
double Ask(double x);
double Answer(double ans);
\end{verbatim}

\section{Dữ liệu và ví dụ gọi hàm}

Nếu miền trong bài là \([0,2]\) và
\[
  f(x)=-(x-1)^2+10,
\]
một chuỗi lời gọi có thể đạt điểm tối đa như sau:

\[
\begin{array}{c|c|c}
\text{Pascal} & \text{C++} & \text{Ý nghĩa} \\
\hline
\texttt{Start;} & \texttt{Start();} & \text{trả về }2,\text{ độ dài miền} \\
\texttt{Ask(0);} & \texttt{Ask(0);} & \text{trả về }9,\ f(0)=9 \\
\texttt{Ask(1.5);} & \texttt{Ask(1.5);} & \text{trả về }9.75 \\
\texttt{Ask(1);} & \texttt{Ask(1);} & \text{trả về }10 \\
\texttt{Ask(2);} & \texttt{Ask(2);} & \text{trả về }9 \\
\texttt{Answer(1);} & \texttt{Answer(1);} & \text{trả đáp án và kết thúc} \\
\end{array}
\]

Ví dụ trên chỉ minh họa cách dùng thư viện, không có ý nghĩa về mặt thuật
toán.

Trong dữ liệu chính thức, \(L\le 10^4\). Đáp án cuối cùng phải có sai số theo
hoành độ so với đỉnh cao nhất không vượt quá \(10^{-3}\). Độ cao của núi nằm
trong đoạn \([0,8848.13]\).

\section{Cách kiểm tra chương trình}

Trong thư mục hiện tại, tạo tệp \texttt{explorer.in}. Tệp này chứa một số
thực \(L\), biểu thị miền xác định của \(f(x)\) là \([0,L]\). Thư viện tương
tác sẽ tự động sinh giá trị của hàm \(f(x)\).

Sau khi chạy chương trình, thư viện ghi toàn bộ cuộc đối thoại vào tệp
\texttt{explorer.log}.

\section{Thông báo lỗi}

Nếu tệp \texttt{explorer.log} chứa một hoặc nhiều dòng bắt đầu bằng
\texttt{Error:}, chương trình đã gặp lỗi trong quá trình chạy. Các lỗi có thể
gồm:
\begin{description}
  \item[\texttt{Error: Input file explorer.in not found}] Không tìm thấy tệp
  \texttt{explorer.in} trong thư mục hiện tại.
  \item[\texttt{Error: Start must call first!}] Chương trình chưa khởi tạo thư
  viện tương tác.
  \item[\texttt{Error: You have called Start twice}] Chương trình gọi
  \texttt{Start} nhiều lần.
  \item[\texttt{Error: Parameter is out of range}] Tham số của \texttt{Ask}
  nằm ngoài miền cho phép.
  \item[\texttt{Error: You have called Ask more than 60 times}] Chương trình
  gọi \texttt{Ask} quá nhiều lần, vượt quá \(60\) lần.
\end{description}

\end{document}
